%% IROS 2027 submission -- IEEE conference format (ieeeconf.cls from PaperCept).
%% IROS allows 6 pages + 2 additional pages (references included in the limit).
%% Deadline: 1 March 2027, 23:59 Pacific.
\documentclass[letterpaper, 10 pt, conference]{ieeeconf}
%\documentclass[a4paper, 10pt, conference]{ieeeconf}

\IEEEoverridecommandlockouts % needed for \thanks
\overrideIEEEmargins          % needed to meet printer requirements

% Fixes "Error 1010/1000" in PaperCept's PDF check for some pdflatex builds.
\pdfobjcompresslevel=0
\pdfminorversion=4

\usepackage[T1]{fontenc}
\usepackage{times}
\usepackage{amsmath}
\usepackage{amssymb}
\let\proof\relax
\let\endproof\relax
\usepackage{amsthm}
\usepackage{graphicx}
\usepackage{booktabs}
\usepackage{multirow}
\usepackage{xcolor}
\usepackage{cite}
\usepackage{url}
\usepackage{tikz}
\usetikzlibrary{arrows.meta, positioning, fit, backgrounds, calc}
\usepackage{algorithm}
\usepackage{algpseudocode}
\usepackage[hidelinks]{hyperref}

% Draft notes -- set \drafttrue to \draftfalse to hide them.
\newif\ifdraft
\drafttrue
\ifdraft
  \newcommand{\todo}[1]{{\color{red}[TODO: #1]}}
  \newcommand{\note}[2]{{\color{blue}[#1: #2]}}
\else
  \newcommand{\todo}[1]{}
  \newcommand{\note}[2]{}
\fi

\newcommand{\robocasa}{RoboCasa}
\newcommand{\segmind}{SegMind}
\newcommand{\eql}{EQL}
\newcommand{\sdt}{SDT}
\newcommand{\pr}{\mathrm{P}}
\newcommand{\dop}{\mathrm{do}}
\newcommand{\tbd}{\textcolor{red}{--}}

\newtheorem{definition}{Definition}

\title{\LARGE \bf
What Went Wrong and How to Fix It:\\
Event-Grounded Causal Circuits for Explainable Robot Learning in Kitchens
}

\author{Sorin Arion$^{1}$, \todo{co-authors}, and Michael Beetz$^{1}$%
\thanks{$^{1}$The authors are with the AICOR Institute for Artificial Intelligence,
        University of Bremen, Bremen, Germany
        {\tt\small \{sorin, beetz\}@uni-bremen.de}}%
}

\begin{document}

\maketitle
\thispagestyle{empty}
\pagestyle{empty}

%%%%%%%%%%%%%%%%%%%%%%%%%%%%%%%%%%%%%%%%%%%%%%%%%%%%%%%%%%%%%%%%%%%%%%%%%%%%%%%%
\begin{abstract}
Robots that learn from their own experience rarely say why an action failed or
what they would change, and robots that explain their failures rarely improve
from them. We present a system that does both on \robocasa{}'s household tasks.
We use \segmind{}, an event detection framework, to segment each episode into
symbolic events over a semantic digital twin: contact, holding, lifting,
placing, containment, opening. \segmind{} annotates the recorded episode with
these events and decides whether the task succeeded, using the same predicates
as \robocasa{}'s own success checks. Interventional causal circuits are fitted
to the robot's annotated runs. They are tractable probabilistic models that
answer $\pr(\text{success} \mid \dop(\cdot))$ exactly. When \segmind{} detects a
failure, it names the failed action and the moment it failed. The circuits then
decide whether the cause was one of the robot's own action parameters or
something in the scene. For a parameter cause they propose corrective
parameters, and for a scene cause they trigger a plan repair. Every decision comes with its explanation: the inference
trace behind the detected event, and the interventional probabilities behind the
blamed cause. Every run is stored in episodic memory in the twin's own classes
and refits the circuits, so the robot learns from its own runs without
demonstrations. We evaluate on the PR2, the Toyota HSR and the PAL TIAGo in
\robocasa{} kitchens. \todo{Headline numbers: annotation accuracy against
\robocasa{}'s success checks, recovery rate against blind resampling, attempts
saved.}
\end{abstract}

%%%%%%%%%%%%%%%%%%%%%%%%%%%%%%%%%%%%%%%%%%%%%%%%%%%%%%%%%%%%%%%%%%%%%%%%%%%%%%%%
\section{Introduction}
\label{sec:introduction}

A robot that puts a mug into a microwave can fail in many ways. It
may stand too far from the counter, grasp the mug too low, or let it slip in
transport. A bowl next to the mug may also block the fingers. To recover, the
robot needs three answers: \emph{what} went wrong, \emph{why} it went wrong, and
\emph{what to change}. To improve, it must keep those answers and learn from
them.

Learned policies trained on large demonstration datasets, the standard approach
on benchmarks such as \robocasa{}~\cite{nasiriany2024robocasa}, give none of
these answers explicitly. A failure is visible only in the final success flag.
Failure-explanation systems based on large vision-language
models~\cite{liu2023reflect,duan2024aha} describe failures in natural language.
Their explanations, however, are not grounded in a model the robot acts on, and
they do not say which parameter to change by how much. Causal approaches to
robot failure~\cite{diehl2022why} identify causes from simulated interventions,
but over a fixed set of variables and for a single action.

We argue that a robot needs two capabilities that share one model of the world.
The first is \emph{event understanding}, which turns what happened during a run
into the symbolic events that define the task, and the second is \emph{causal
reasoning}, which links those events to the robot's own choices and to the
objects in the scene. We build on three of our earlier works. Each of them
solves one part of this problem, but none of them solves it alone.
\segmind{}~\cite{arion2026segmind} segments a run into explained events over a
semantic digital twin (\sdt{})~\cite{kuempel2021sdt}, but only says \emph{what}
happened. Interventional causal circuits~\cite{vasantakumaar2026icc} turn a
failed action into the parameter that caused it and a corrective value. They
need an external pass/fail tester and a fixed set of variables, and were shown
on one pick-and-place. Relational causal circuits~\cite{anon2027relational} lift
these circuits to scenes with a varying number of objects, so that a cause can
be ``how crowded the target is'' or ``the neighbour on the fingers' closing
axis''.

\robocasa{} needs all three. Its success checks are written in terms of the same
relations that \segmind{} detects, for example whether an object lies inside a
receptacle or whether a door has been opened far enough. Its tasks
are sequential, and its scenes are relational: \todo{87} of its task files place
a varying number of distractor objects on counters, in cabinets, in the sink and
in the fridge.

\paragraph*{Contributions}
\begin{enumerate}
  \item \textbf{Event-based annotation and scoring of robot episodes.}
        \segmind{} annotates recorded episodes with typed, explained events over
        the \sdt{}. The same events decide task success, using the predicates of
        \robocasa{}'s own success checks. They also act as the tester the causal
        circuits learn from, and fill in what happened to every other object in
        the scene.
  \item \textbf{Learnable and explainable causal diagnosis at two levels.}
        Flat causal circuits over each action type's parameters, and relational
        causal circuits grounded for the objects actually around the target,
        decide whether the robot's parameters or the scene caused a failure. A
        parameter cause is answered with a one-shot correction, a scene cause
        with a plan repair. Both are refitted from the robot's own runs.
  \item \textbf{One object model from execution to causal query.} Success
        predicates, expected action effects, episodic memory and causal queries
        are all expressed in the Entity Query Language (\eql{})~\cite{bassiouny2026krrood}
        over the twin's classes. No hand-built feature table is needed.
\end{enumerate}

We evaluate the system on three mobile manipulators, the PR2, the Toyota HSR
and the PAL TIAGo, on \robocasa{} tasks in \robocasa{}'s kitchens as built.
\todo{Task count; real-robot runs if any.}

%%%%%%%%%%%%%%%%%%%%%%%%%%%%%%%%%%%%%%%%%%%%%%%%%%%%%%%%%%%%%%%%%%%%%%%%%%%%%%%%
\section{Related Work}
\label{sec:related-work}

\paragraph*{Household manipulation benchmarks}
\robocasa{}~\cite{nasiriany2024robocasa} builds on
robosuite~\cite{zhu2020robosuite} and MuJoCo~\cite{todorov2012mujoco}. It
provides procedurally composed kitchens, atomic and composite tasks, and large
demonstration datasets for policy learning. Success is decided by hand-written
checks inside each task class. We keep \robocasa{}'s kitchens and success
thresholds, but replace the end-to-end policy with knowledge-based
plans~\cite{beetz2010cram} whose continuous parameters are learned. We also
make the success checks explicit and inspectable as \segmind{} events.

\paragraph*{Event segmentation and episode annotation}
Event segmentation theory~\cite{zacks2007event,kurby2008segmentation} holds that
people parse continuous activity into discrete events at points of significant
change. In robotics, rule-based, learned and foundation-model-based segmenters
annotate demonstrations and robot runs. Learned segmenters seldom expose the
reasoning behind a boundary. \segmind{}~\cite{arion2026segmind} detects events
with statecharts~\cite{stelter2025motionstatecharts} whose detectors are \eql{}
rules over the \sdt{}, so every event carries an inference trace.

\paragraph*{Failure explanation and recovery}
REFLECT~\cite{liu2023reflect} summarises multisensory robot experience
hierarchically and uses a language model to explain and correct failures.
AHA~\cite{duan2024aha} fine-tunes a vision-language model to detect and reason
about manipulation failures. Diehl and Ramirez-Amaro~\cite{diehl2022why} learn a
causal Bayesian network from simulated interventions to explain failures in
stacking. Our explanations are tied to the model the robot acts on. The events
come from the twin, and the causes are variables of the robot's plan or of the
scene, with an exact interventional probability for each.

\paragraph*{Tractable causal inference}
Probabilistic circuits~\cite{choi2020probabilistic} answer marginal and
conditional queries in time linear in their size. Under marginal
determinism~\cite{wang2023compositional}, they also answer backdoor-adjusted
interventional queries exactly. Interventional causal
circuits~\cite{vasantakumaar2026icc} learn such circuits as joint probability
trees~\cite{nyga2023jpt} from successful robot runs and use them to diagnose and
correct failed plans. Relational causal circuits~\cite{anon2027relational}
ground a circuit learned over a class and its collections for the objects a
query names, building on relational sum-product
networks~\cite{nath2015relational}.

%%%%%%%%%%%%%%%%%%%%%%%%%%%%%%%%%%%%%%%%%%%%%%%%%%%%%%%%%%%%%%%%%%%%%%%%%%%%%%%%
\section{Approach}
\label{sec:approach}

\begin{figure*}[t]
  \centering
  \begin{tikzpicture}[
      node distance=4mm and 7.5mm,
      box/.style={draw, rounded corners=2pt, align=center, font=\footnotesize,
                  minimum height=10mm, text width=20.5mm, fill=white},
      mem/.style={box, fill=black!6},
      arr/.style={-{Stealth[length=2mm]}, thick},
      lab/.style={font=\scriptsize\itshape, align=center}]
    \node[box] (propose) {\textbf{Propose}\\parameters from\\flat circuits};
    \node[box, right=of propose] (execute) {\textbf{Execute}\\CRAM plan in\\MuJoCo / robot};
    \node[box, right=of execute] (segment) {\textbf{Segment}\\\segmind{} events\\over the \sdt{}};
    \node[box, right=of segment] (localise) {\textbf{Localise}\\first missing or\\unexpected event};
    \node[box, right=of localise] (diagnose) {\textbf{Diagnose}\\flat + relational\\causal circuits};
    \node[box, right=of diagnose] (correct) {\textbf{Correct / repair}\\region or\\plan repair};
    \node[mem, below=10mm of segment, text width=60mm] (memory)
      {\textbf{Episodic memory} (ORM): attempts, parameters, neighbours,
       annotated events, all as objects of the twin's classes};
    \draw[arr] (propose) -- (execute);
    \draw[arr] (execute) -- (segment);
    \draw[arr] (segment) -- node[lab, above] {failure} (localise);
    \draw[arr] (localise) -- (diagnose);
    \draw[arr] (diagnose) -- (correct);
    \draw[arr] (correct.north) -- ++(0,4mm) -| node[lab, above, pos=0.25] {retry with corrected parameters or repaired plan} (propose.north);
    \draw[arr] (segment) -- node[lab, right] {annotate} (memory);
    \draw[arr] (memory.west) -| node[lab, below, pos=0.25] {refit (learn)} (propose.south);
    \draw[arr] (memory.east) -| node[lab, below, pos=0.25] {fit / ground} (diagnose.south);
  \end{tikzpicture}
  \caption{The loop. A CRAM plan proposes each action's parameters from that
  action's learned flat circuit and executes them. \segmind{} annotates the run with
  events over the semantic digital twin and scores it against \robocasa{}'s
  success predicates. On failure, the first missing or unexpected event
  localises the failed action. Flat and relational causal circuits then decide
  whether the robot's parameters or the scene caused it, and the plan is
  corrected or repaired. Every run, annotated, goes into episodic memory and
  refits the models.}
  \label{fig:overview}
\end{figure*}

Fig.~\ref{fig:overview} shows the loop; Alg.~\ref{alg:loop} states it. We first
describe the shared world model (Sec.~\ref{sec:sdt}), then event-based
annotation (Sec.~\ref{sec:segmind}), the causal circuits
(Sec.~\ref{sec:circuits}), and how they are combined for diagnosis, correction
and learning (Sec.~\ref{sec:loop}).

\subsection{One Object Model: Twin, Tasks and Memory}
\label{sec:sdt}

The semantic digital twin represents the scene as a world
$W = (O, G, R)$~\cite{kuempel2021sdt}. $O$ is a set of bodies, $G$ assigns each
body a geometry and a 6-DoF pose, and $R$ is a kinematic graph whose
connections carry joint states. Bodies carry semantic annotations such as
\texttt{Drawer}, \texttt{Fridge}, \texttt{Handle} and \texttt{CounterTop}. We
load \robocasa{} kitchens into the twin directly from \robocasa{}'s MJCF and
fixture metadata, so robot, kitchen and task objects share one model.
\todo{Loader details: layouts/styles, mass and joint-limit fixes.}

All knowledge is written in \eql{}~\cite{bassiouny2026krrood}, a query language
embedded in Python that ranges over the twin's classes directly. \robocasa{}'s
success checks are restated as \eql{} conditions over \segmind{}'s predicates.
For example, the
condition of \texttt{PickPlaceCounterToSink} is that the object is contained in
the sink basin and no longer in contact with the gripper. Episodic memory stores
every run through an object-relational mapping (ORM) of the same classes: the action,
its parameters, the target, the target's neighbours and the events. A
relational circuit can therefore be fitted on memory without flattening it into
a table.

\subsection{\segmind{}: Event-Based Annotation of Episodes}
\label{sec:segmind}

\segmind{}~\cite{arion2026segmind} consumes a stream of twin states, either live
during execution or replayed from a recorded episode. It produces a log of typed
events, each with a start and end time, the objects involved, and the \eql{}
inference trace that produced it. Detectors are organised as event detector
statecharts~\cite{stelter2025motionstatecharts}. A detector becomes active when
its start condition, a formula over other detectors' states, holds, so the
detection of composite events needs no hand-written control flow. There are
three levels:
\begin{itemize}
  \item \emph{atomic} events: contact and loss of contact, start and stop of
        translation and rotation;
  \item \emph{spatial-relation} events: support, containment and their loss,
        and the opening and closing of articulated fixtures;
  \item \emph{coarse} events composed of these: pick-up (translation plus loss
        of support), placing (stop of motion plus support), insertion.
\end{itemize}

\paragraph*{Annotating episodes}
Applied to a recorded episode, the event log is a temporal annotation of that
episode. It contains which object was picked up when, where it was placed, which
door was opened, and which other objects were touched or moved on the way.
Because detection runs on the twin and not on pixels, the same detectors
annotate any robot and any episode source: our own PR2, HSR and TIAGo runs, and
\todo{\robocasa{}'s demonstration datasets, replayed into the twin}.

\paragraph*{Expected effects and success}
Each action type declares the events it is expected to cause.
\texttt{PickUpAction} expects holding and lifting of its object;
\texttt{PlaceAction} expects a placing on its target; \texttt{OpenAction}
expects the fixture's opening past a threshold. A task succeeds when its \eql{}
success condition holds over the annotated events at \robocasa{}'s own
thresholds. Thus one definition of success scores the benchmark, rejects a run
and serves as the effect variable of the causal model. A score, a rejection and
a diagnosis can never disagree about what success meant.

\subsection{Interventional Causal Circuits}
\label{sec:circuits}

\paragraph*{Flat circuits over action parameters}
For each action type $a$ (pick up, place, open, close, \dots) we fit a joint
probability tree (JPT)~\cite{nyga2023jpt}, a probabilistic circuit, over the action's
continuous parameters $X^a = (X_1, \dots, X_n)$ and the outcome variables $E$. For a pick-up, the parameters are the
standing pose, approach direction, grasp height and arm. The outcomes are
\segmind{}'s events, e.g.\ whether holding and lifting were detected. The
circuit is structured by a marginal determinism tree derived from the causal
graph: every sum unit at a level is support-deterministic over that level's
query variables, i.e.\ its children's supports over those variables are
disjoint~\cite{wang2023compositional}. We verify this property on every fitted
circuit before querying it. Under it, the backdoor adjustment
\begin{equation}
  \pr(E \mid \dop(X_i = x)) = \sum_{z} \pr(E \mid X_i = x, Z = z)\,\pr(Z = z)
  \label{eq:backdoor}
\end{equation}
for an adjustment set $Z$~\cite{pearl2009causality} is computed exactly, in time polynomial in the
circuit's size. The result is itself a circuit over $(X_i, E)$. For randomised
training runs $Z = \emptyset$ and \eqref{eq:backdoor} reduces to conditioning.
For deployment data where scene state confounds parameter choice, $Z$ contains
that scene state.

\paragraph*{Relational circuits over the scene}
A flat circuit can only blame the robot's own parameters. However, a correction
to those parameters cannot help when a bowl stands in the way. We therefore
also fit a relational causal circuit~\cite{anon2027relational} over the attempt
class and its collection of neighbours. Each neighbour has its side of the
fingers' closing axis, a distance band, and whether \segmind{} saw it moved or
touched. The circuit is learned once over the class. For a given failed
attempt, it is grounded for exactly the neighbours present. A scene cause can
then be an aggregate (``the number of objects within 10\,cm of the target'') or
one part (``the neighbour on the closing axis''). The circuit keeps the guarantees
\eqref{eq:backdoor} needs.

\paragraph*{Valid adjustment from the timeline}
Backdoor adjustment must not adjust for a descendant of the cause. \segmind{}'s
timestamps make the split explicit. Scene state recorded before an action starts
cannot be a descendant of that action's parameters and is admissible in $Z$.
Anything after the action starts is not.

\subsection{Diagnosis, Correction and Learning}
\label{sec:loop}

\begin{algorithm}[t]
  \caption{Event-grounded causal recovery}
  \label{alg:loop}
  \small
  \begin{algorithmic}[1]
    \Require task $\tau$, plan $\pi_\tau$, circuits $\{\mathcal{C}_a\}$, relational circuit $\mathcal{R}$, memory $\mathcal{M}$
    \Loop
      \State $\theta \sim \mathcal{C}_a(\cdot \mid E = \text{success})$ for each action $a \in \pi_\tau$ \Comment{propose}
      \State $\xi \gets \textsc{Execute}(\pi_\tau, \theta)$
      \State $\mathcal{E} \gets \textsc{SegMind}(\xi)$ \Comment{annotate}
      \State $\mathcal{M} \gets \mathcal{M} \cup \{(\theta, \xi, \mathcal{E})\}$
      \If{$\textsc{Success}_\tau(\mathcal{E})$} \Return success \EndIf
      \State $(a, t) \gets$ first missing or unexpected event in $\mathcal{E}$ \Comment{localise}
      \State $(i^\ast, q_{i^\ast}, \rho) \gets \textsc{Diagnose}(\mathcal{C}_a, \theta_a)$
      \State $(s^\ast, q_{s^\ast}) \gets \textsc{Diagnose}(\mathcal{R}.\textsc{ground}(\text{scene at } t))$
      \If{$q_{s^\ast} < q_{i^\ast}$} $\pi_\tau \gets \textsc{Repair}(\pi_\tau, s^\ast)$
      \Else{} $\theta_{a,i^\ast} \gets$ midpoint of $\rho$ \EndIf
      \State refit $\mathcal{C}_a$, $\mathcal{R}$ on $\mathcal{M}$ \Comment{learn}
    \EndLoop
  \end{algorithmic}
\end{algorithm}

\paragraph*{Localise}
The first expected event that is missing, or the first unexpected one, names
the failed action $a$ and the time $t$ of failure. Examples are a pick-up
without lifting, or a loss of holding before the placing. The event's \eql{}
inference trace states which rule fired with which bindings.

\paragraph*{Diagnose}
For each parameter $X_i$ of $a$ with observed value $x_i$, we query the
interventional circuit of $X_i$ on the effect at a small interval around $x_i$:
\begin{equation}
  q_i = \pr\big(X_i \in [x_i - \epsilon, x_i + \epsilon]\big)
  \text{ in } \pr(X_i, E \mid \dop(X_i)).
  \label{eq:diagnosis}
\end{equation}
Here $q_i = 0$ means that $x_i$ lies outside the support of successful
training runs. The primary cause is $i^\ast = \arg\min_i q_i$. The recommended
region $\rho$ is the support region of $X_{i^\ast}$ with the highest
interventional probability. The relational circuit is queried in the same way
for the scene variables, grounded at the scene before $t$. Values outside the
training support are flagged rather than corrected.

\paragraph*{Correct or repair}
A parameter cause gets a one-shot correction: $X_{i^\ast}$ is set inside
$\rho$, and the action is retried. A scene cause gets a plan repair from CRAM:
set the offending neighbour aside, approach from the other side, or use the
other arm.

\paragraph*{Explain}
The explanation of a failure is the pair of the event trace (\emph{what} failed,
\emph{when}, and by which rule) and the causal report (\emph{which} variable,
its observed value, $q_i$ for every candidate, and the recommended region).
Both are produced by the models the robot acts on, not generated after the
fact. \todo{Fig.: an example explanation from a real run.}

\paragraph*{Learn}
Every run, successful or not, is stored annotated in episodic memory. Successful
action instances refit the per-action JPTs that the plan samples from, and all
instances refit the circuits. Circuits are fitted per action type, not per
task. Relational circuits are grounded per scene. A handful of models therefore
covers every \robocasa{} task built from these actions, whatever objects stand
around. The robot learns without demonstrations.

%%%%%%%%%%%%%%%%%%%%%%%%%%%%%%%%%%%%%%%%%%%%%%%%%%%%%%%%%%%%%%%%%%%%%%%%%%%%%%%%
\section{Experiments}
\label{sec:experiments}

We evaluate four questions.
\textbf{Q1}: Does \segmind{} understand the scene and the task well enough to
annotate episodes and score them like \robocasa{} does?
\textbf{Q2}: Do the causal circuits find the cause of a failure and fix it in
fewer attempts than uninformed retries?
\textbf{Q3}: Do relational circuits separate scene causes from parameter causes?
\textbf{Q4}: Does the robot improve from its own runs, and do the explanations
stay correct as it learns?

\subsection{Setup}
\label{sec:setup}

\paragraph*{Robots}
We use three mobile manipulators with different kinematics and grippers: the
PR2 (two 7-DoF arms, parallel grippers), the Toyota HSR (one 5-DoF arm on an
omnidirectional base with a telescopic torso), and the PAL TIAGo (one 7-DoF arm,
lifting torso). All three run the same CRAM plans. Motions are generated by
\todo{Giskard} and base placement by graphs of convex
sets~\cite{marcucci2023gcs}. \todo{TIAGo gripper variant.}

\paragraph*{Environment and tasks}
All runs are in MuJoCo in \robocasa{}'s kitchens as built (0.92\,m counters),
loaded into the twin. Lowering a kitchen so that a robot can reach changes the
benchmark; such runs are reported separately as a labelled variant. We use
\robocasa{} atomic tasks from the pick-and-place, drawer and door families,
listed in Table~\ref{tab:tasks}, \todo{and composite tasks}. Perception is
\todo{known state / simulated visibility / RoboKudo~\cite{mania2024robokudo} on
rendered images}.

\paragraph*{Data}
\todo{Number of episodes per robot and task; random parameter sampling for
training; fault injection (pose noise, object mass, friction, stiff joints) for
the test set; layouts and styles held out.}

\begin{table}[t]
  \centering
  \caption{\robocasa{} tasks per robot in \robocasa{}'s kitchens as built:
  success rate over \todo{$N$} episodes, as scored by \segmind{}
  \todo{fill in}.}
  \label{tab:tasks}
  \setlength{\tabcolsep}{4pt}
  \begin{tabular}{llccc}
    \toprule
    Family & Task & PR2 & HSR & TIAGo \\
    \midrule
    \multirow{2}{*}{Pick \& place}
      & PickPlaceCounterToSink    & \tbd & \tbd & \tbd \\
      & PickPlaceCounterToCabinet & \tbd & \tbd & \tbd \\
    \midrule
    \multirow{2}{*}{Drawer}
      & OpenDrawer  & \tbd & \tbd & \tbd \\
      & CloseDrawer & \tbd & \tbd & \tbd \\
    \midrule
    \multirow{3}{*}{Door}
      & OpenMicrowave  & \tbd & \tbd & \tbd \\
      & CloseMicrowave & \tbd & \tbd & \tbd \\
      & CloseOven      & \tbd & \tbd & \tbd \\
    \bottomrule
  \end{tabular}
\end{table}

\subsection{Q1: Scene and Task Understanding with \segmind{}}
\label{sec:exp-segmind}

\paragraph*{Protocol}
We replay every recorded episode into \segmind{} and compare its annotations
against two references. The first is \robocasa{}'s own \texttt{\_check\_success}
evaluated in the same simulation, which measures agreement on task success. The
second is event ground truth from the simulator's contact and joint states, which
gives precision, recall and temporal error of event boundaries. \todo{Also
annotate \robocasa{}'s demonstration datasets.}

\paragraph*{Results}
\todo{Table~\ref{tab:segmind}; discuss disagreements, especially the placing
pairing and race conditions fixed for this paper.}

\begin{table}[t]
  \centering
  \caption{Episode annotation by \segmind{}: agreement with \robocasa{}'s
  success checks, and event detection against simulator ground truth
  \todo{fill in}.}
  \label{tab:segmind}
  \setlength{\tabcolsep}{3.5pt}
  \begin{tabular}{lccccc}
    \toprule
    Robot & Episodes & Success agr. & Event P & Event R & $\Delta t$ [s] \\
    \midrule
    PR2   & \tbd & \tbd & \tbd & \tbd & \tbd \\
    HSR   & \tbd & \tbd & \tbd & \tbd & \tbd \\
    TIAGo & \tbd & \tbd & \tbd & \tbd & \tbd \\
    \bottomrule
  \end{tabular}
\end{table}

\subsection{Q2: Diagnosing and Correcting Parameter Failures}
\label{sec:exp-recovery}

\paragraph*{Protocol}
We inject parameter faults into otherwise successful plans and compare three
recovery strategies, each allowed up to \todo{$k$} attempts.
\textbf{Blind resampling} draws new parameters from the JPT, the baseline of
\cite{vasantakumaar2026icc}. \textbf{Circuit + final state} is the flat causal
circuit with a final-state tester (object height only), as
in~\cite{vasantakumaar2026icc}. \textbf{Circuit + \segmind{}} is our system.
We report whether the injected parameter is named as the primary cause, the
recovery rate, and the mean number of attempts. \todo{Replanning from scratch
as a further baseline.}

\paragraph*{Results}
\todo{Table~\ref{tab:recovery}.}

\begin{table}[t]
  \centering
  \caption{Recovery from injected parameter faults, averaged over tasks
  \todo{fill in}. Cause: injected parameter named as primary cause.}
  \label{tab:recovery}
  \setlength{\tabcolsep}{3pt}
  \begin{tabular}{llccc}
    \toprule
    Robot & Method & Cause [\%] & Recovered [\%] & Attempts \\
    \midrule
    \multirow{3}{*}{PR2}
      & Blind resampling         & --   & \tbd & \tbd \\
      & Circuit + final state    & \tbd & \tbd & \tbd \\
      & Circuit + \segmind{}     & \tbd & \tbd & \tbd \\
    \midrule
    \multirow{3}{*}{HSR}
      & Blind resampling         & --   & \tbd & \tbd \\
      & Circuit + final state    & \tbd & \tbd & \tbd \\
      & Circuit + \segmind{}     & \tbd & \tbd & \tbd \\
    \midrule
    \multirow{3}{*}{TIAGo}
      & Blind resampling         & --   & \tbd & \tbd \\
      & Circuit + final state    & \tbd & \tbd & \tbd \\
      & Circuit + \segmind{}     & \tbd & \tbd & \tbd \\
    \bottomrule
  \end{tabular}
\end{table}

\subsection{Q3: Robot or Scene?}
\label{sec:exp-relational}

\paragraph*{Protocol}
We populate \robocasa{} counters with $0$ to \todo{$N$} distractors around the
target, which produces failures with scene causes (blocked fingers, disturbed
neighbours). We compare a flat circuit on a hand-flattened table (counts only,
or one block of columns per neighbour) against the grounded relational circuit,
on the same runs. We measure whether the cause is attributed correctly to the
parameters or to the scene. We also compare the success of plan repair against
parameter correction. \todo{Deep-set g-computation as in~\cite{anon2027relational}.}

\paragraph*{Results}
\todo{Confusion matrix of attributed cause type; recovery per distractor count.}

\subsection{Q4: Learning and Explainability}
\label{sec:exp-learning}

\paragraph*{Protocol}
Starting from \todo{$n_0$} random-parameter runs, the robot runs tasks
repeatedly, storing every annotated run and refitting the models after every
\todo{$m$} episodes. We report task success and the accuracy of cause
attribution against the number of episodes in memory. For explainability, we
check that each explanation is \emph{faithful}: applying the recommended
correction to the failed run's parameters, and nothing else, makes the run
succeed in \todo{X}\% of cases. \todo{A user study, or a check of explanations
against injected causes.}

\paragraph*{Results}
\todo{Learning curves per robot (one figure, three panels).}

\subsection{Composite Tasks}
\label{sec:exp-composite}
\todo{Recovery inside composite tasks: failures per step; fixing the failed
step against restarting the task.}

%%%%%%%%%%%%%%%%%%%%%%%%%%%%%%%%%%%%%%%%%%%%%%%%%%%%%%%%%%%%%%%%%%%%%%%%%%%%%%%%
\section{Discussion and Limitations}
\label{sec:discussion}

\emph{Interventions are one cause at a time.} The circuits intervene on one
discrete or interval-valued cause at a time. Relations between neighbours enter
only through their attributes and counts. Joint causes, such as a grasp height
that fails only with a given approach, are not yet diagnosed.

\emph{Out-of-support values are flagged, not fixed.} A failure whose parameters
lie outside every training run gets $q_i = 0$ for every variable. The robot
reports this case rather than guessing.

\emph{The twin can be wrong.} \segmind{} reasons over the twin. With imperfect
perception, a misdetected object pose propagates into the events. We treat the
perceived pose as an observed variable, so misperception can show up as a cause.
\todo{Quantify with perception levels.}

\emph{Reach.} The HSR cannot reach \robocasa{}'s wall cabinets or the inside of
some drawers at 0.92\,m counters. We report these tasks as blocked rather than
change the kitchen.

%%%%%%%%%%%%%%%%%%%%%%%%%%%%%%%%%%%%%%%%%%%%%%%%%%%%%%%%%%%%%%%%%%%%%%%%%%%%%%%%
\section{Conclusion}
\label{sec:conclusion}

We connected event-based annotation of robot episodes with interventional
causal circuits. On \robocasa{}'s tasks, a PR2, an HSR and a TIAGo detect their
failures, localise them to an action and a moment, decide whether their own
parameters or the scene caused them, and correct or repair the plan. Each step
is explained by the models the robot acts on, and the robot learns from its own
runs without demonstrations. \todo{Summary of results.} Future work covers
joint causes, real-robot runs in the lab kitchen, and annotating large
demonstration datasets to fit circuits for learned policies.

%%%%%%%%%%%%%%%%%%%%%%%%%%%%%%%%%%%%%%%%%%%%%%%%%%%%%%%%%%%%%%%%%%%%%%%%%%%%%%%%
\section*{Acknowledgment}
\todo{Funding.}

% Balance the columns on the last page if needed:
%\addtolength{\textheight}{-3cm}

\bibliographystyle{IEEEtran}
\bibliography{references}

\end{document}
